\section{Introduction}
{ 
Quantum spin liquids (QSLs) represent quantum phases beyond the conventional Ginzburg--Landau paradigm, characterized not by local order parameters but by long-range quantum entanglement, fractionalized excitations, and emergent gauge structures~\cite{Savary-Balent_2017,Zhou-Kanoda-Ng_2017}. Unlike conventional ordered phases, whose essential properties can often be captured through spontaneous symmetry breaking, the defining characteristics of a QSL are encoded in the structure of its many-body quantum state. Consequently, much of our theoretical understanding of QSLs has been developed through their zero-temperature ground-state wave functions and the genuinely quantum structures contained within them.

Quantum entanglement provides one of the most powerful frameworks for extracting such structures~\cite{ZengChenZhouWen2019,Wen_2017,Laflorencie_2016,Amico-Fazio-Osterloh-Vedral_2008,Eisert-Cramer-Plenio_2010,Lyu-Chandorkar-Kapoor-Takei-Sorensen-Witczak-Krempa_2026}. For a pure ground state, the reduced density matrix (RDM) of a subsystem provides a complete characterization of its bipartite entanglement. Its eigenvalues determine the entanglement entropy (EE) and entanglement spectrum (ES), which have become central tools for identifying topological order~\cite{Levin-Wen_2006,Kitaev-Preskill_2006}, bulk--boundary correspondence~\cite{Li-Haldane,Li-Haldane-1,Li-Haldane-2}, and other manifestations of long-range quantum entanglement. This ground-state and entanglement-based perspective has led to remarkable progress in understanding QSLs at zero temperature.

At the same time, the success of this perspective also exposes a limitation of our present theoretical understanding rather than a fundamental limitation of the underlying physics. The characteristic quantum structures of QSLs are presently understood most sharply for pure ground states, where entanglement admits particularly direct characterizations. Once the state becomes mixed, as at finite temperature or in an open quantum system, several inequivalent measures of quantum and mixed-state entanglement are possible, but no single quantity plays the same general role as the bipartite EE of a pure state. This raises a broader question: which characteristic structures of a QSL can be identified directly at the level of its physical density matrix, which remains well defined for pure, mixed, thermal, and open-system states? A necessary first step toward addressing this broader question is to understand how the density matrix is organized already in an exactly solvable QSL at zero temperature.

This question also has an important operational aspect. While the density matrix specifies a quantum state, its expression in the physical operator algebra connects that state to physically accessible observables and operations. For a quantum spin system, these are naturally formulated in terms of the physical spin degrees of freedom. An explicit density matrix expressed in spin operators therefore provides a natural interface with quantum-information and quantum-resource frameworks, in which quantum properties are characterized relative to physically meaningful operations~\cite{Nielsen_Chuang_2010_sec8,Chitambar-Gour,Leone-Oliviero-Hamma_2022}. Thus, understanding a QSL at the level of its physical spin density matrix is relevant not only for extending its characterization beyond pure ground states, but also for identifying quantum structures that are operationally accessible in the underlying spin system.

The Kitaev honeycomb model~\cite{Kitaev_2006} provides an ideal setting in which to establish such a physical density-matrix description. As one of the few exactly solvable interacting spin models realizing a $\mathbb{Z}_2$ QSL, it has become a paradigmatic platform for understanding fractionalization, emergent gauge fields, topological order, and anyonic excitations~\cite{Savary-Balent_2017,Zhou-Kanoda-Ng_2017}, with growing connections to candidate Kitaev materials and experiments~\cite{Matsuda-Shibauchi-Kee_2025,Kasahara-etal_2018,Sandilands-Tian-Plumb-Kim-Burch_2015,Nasu-Knolle-Kovrizhin-Motome-Moessner_2016}. Through the Majorana-fermion representation, the interacting spin problem is mapped to free Majorana fermions coupled to a static $\mathbb{Z}_2$ gauge field, providing a powerful framework for calculating its ground-state entanglement properties. In particular, the EE was shown to separate exactly into matter and gauge contributions~\cite{Yao-Qi_2010}, while the ES can be obtained from fermionic correlation functions using standard free-fermion techniques~\cite{Jin-Korepin_2004,Vidal-Latorre-Rico-Kitaev_2003,Latorre-Rico-Vidal_2004,Peschel_2003,Peschel-Eisler_2009,Chung-Peschel_2001,Cheong-Henley_2004}. These developments establish in considerable detail \emph{what} the entanglement properties of the Kitaev spin liquid are.

What remains less transparent is the physical operator organization underlying these results. The conventional entanglement construction is formulated in an enlarged Majorana Hilbert space with gauge redundancy, where the physical state is recovered after imposing gauge constraints and the gauge degrees of freedom intersecting the entanglement boundary require a special treatment~\cite{Yao-Qi_2010}. Consequently, although the EE and ES can be calculated exactly, their underlying structure is not manifest directly in the original spin Hilbert space. In particular, it remains desirable to understand, within a unified physical-spin formulation, how the density matrix is organized, what microscopic structure produces the matter--gauge separation of the EE, why the RDM develops its characteristic block structure and degeneracies, and how gauge constraints and boundary degrees of freedom are encoded directly in the physical density matrix. The recently emphasized role of $\mathbb{Z}_2$ one-form Wilson symmetry~\cite{Xu-Rakovszky-Knap-Pollmann_2025} further suggests that these structures may admit a direct symmetry-based interpretation. These questions concern not the computability of the known entanglement quantities, but the physical organization of the quantum state from which they arise.

In this work, we answer these questions by deriving the exact DM of the Kitaev quantum spin liquid entirely within the original spin Hilbert space. {We show that the DM, as well as the RDM, expressed in terms of spin-string operators naturally implies the separation between the matter and gauge sectors in the entanglement structure. Furthermore, we demonstrate that the RDM is organized into diagonal blocks arising from boundary degrees of freedom constrained jointly by Gauss's law and the exact $\mathbb{Z}_2$ 1-form Wilson symmetry recently discussed in Ref.~\cite{Xu-Rakovszky-Knap-Pollmann_2025}. This structure provides a microscopic origin of the characteristic degeneracy of the entanglement spectrum, identifying it directly with the gauge-constrained boundary degrees of freedom of the physical spin system. In particular, the method of evaluating EE proposed in Ref.~\cite{Yao-Qi_2010} was made possible by redefining the gauge degrees of freedom at the boundary, a technique however seeming applicable only when the subsystem contains an even number of sites since the Hilbert space of an odd number of Majorana fermions would be ill-defined. Our physical spin formulation resolves this ambiguity by treating the matter and gauge sectors on an equal footing. This provides a unified framework applicable to both even and odd subsystem sizes and clarifies the subtle role of subsystem fermion parity in EE that remains hidden in the standard Majorana representation.
}

Beyond providing an exact spin representation of the DM, our work establishes a microscopic organization of the physical quantum state that complements the established Majorana description. More broadly, it identifies the DM itself as a natural object for revealing structural properties of quantum spin liquids. We expect this perspective to provide a useful foundation for investigating which characteristic structures identified at zero temperature survive, evolve, or disappear in finite-temperature and more general mixed-state settings, thereby contributing toward a more unified understanding of quantum spin liquids beyond the pure-ground-state framework.

The remainder of this paper is organized as follows. In Sec.~\ref{sec:model-intro}, we briefly review the Kitaev model and establish the notation. In Sec.~\ref{sec:density-matrix}, we derive the exact DM in the spin representation and demonstrate its organization by equivalence classes of spin-string operators. In Sec.~\ref{sec:entanglement-structures}, we derive the RDM, recover the previously known EE, identify the origin of its block-diagonal structure and entanglement-spectrum degeneracy, and discuss these results from the perspective of higher-form symmetry~\cite{Xu-Rakovszky-Knap-Pollmann_2025}. Finally, Sec.~\ref{sec:summary} summarizes our results and discusses possible future directions.
}